\documentclass[11pt,a4paper]{article}
\usepackage[utf8]{inputenc}
\usepackage[T1]{fontenc}
\usepackage{lmodern}
\usepackage[portuguese]{babel}
\usepackage[a4paper,margin=2.5cm]{geometry}
\usepackage{xcolor}
\definecolor{TextEspresso}{HTML}{4A4238}
\definecolor{NudeTaupe}{HTML}{BFAFA6}
\definecolor{SoftBlush}{HTML}{D9C5B2}
\usepackage{titlesec}
\usepackage{enumitem}
\usepackage{listings}
\usepackage{hyperref}
\hypersetup{colorlinks=true, linkcolor=TextEspresso, urlcolor=TextEspresso}

\titleformat{\section}{\color{TextEspresso}\Large\bfseries}{\thesection}{0.5em}{}[{\color{NudeTaupe}\titlerule[0.8pt]}]
\titleformat{\subsection}{\color{TextEspresso}\large\bfseries}{\thesubsection}{0.5em}{}

\lstset{
  basicstyle=\ttfamily\small,
  breaklines=true,
  frame=single,
  rulecolor=\color{NudeTaupe},
  columns=fullflexible,
  showstringspaces=false
}

\newcommand{\guiaotitle}[2]{%
  \begin{center}
  {\Large\bfseries\color{TextEspresso} #1}\\[0.3cm]
  {\normalsize\color{NudeTaupe} #2}\\[0.2cm]
  {\color{NudeTaupe}\rule{4cm}{0.5pt}}
  \end{center}
  \vspace{0.5cm}
}

\begin{document}
\guiaotitle{Guião 12 -- Monitorização com Net-SNMP}{Infraestruturas de Computação na Nuvem, Aula 12}

\section{Objetivos}
\begin{itemize}[itemsep=0.2cm]
  \item Instalar e configurar o agente SNMP (\texttt{snmpd}) numa VM Linux do Proxmox.
  \item Definir uma \textit{community string} no ficheiro \texttt{snmpd.conf} e reiniciar o serviço.
  \item Consultar OIDs standard a partir de outra VM com \texttt{snmpget} e \texttt{snmpwalk} (uptime do sistema e tabela de interfaces de rede da MIB-II).
  \item Traduzir OIDs numéricos em nomes legíveis e distinguir tipos de dados (Counter vs Gauge).
  \item Configurar o envio de uma trap SNMP simples para um recetor de traps.
  \item Configurar um utilizador SNMPv3 com autenticação e cifra, e comparar o tráfego com o do SNMPv2c.
\end{itemize}

\section{Pré-requisitos e Material}
\begin{itemize}[itemsep=0.2cm]
  \item \textbf{Duas VMs Linux no Proxmox}, na mesma rede: uma no papel de \emph{agente} (a máquina gerida) e outra no papel de \emph{manager}. Pode reutilizar VMs de aulas anteriores.
  \item Acesso SSH a ambas, com privilégios de \texttt{sudo}.
  \item Pacotes \texttt{snmp} e \texttt{snmp-mibs-downloader} na VM manager, e \texttt{snmpd} na VM agente.
  \item Ligação de rede entre as duas VMs, com as portas UDP 161 (SNMP) e UDP 162 (traps) acessíveis.
\end{itemize}

\textbf{Importante:} configure o agente SNMP nas \emph{suas} VMs, nunca no servidor Proxmox propriamente dito. O servidor é partilhado por toda a turma e a sua configuração não deve ser alterada. Se quiser observar dados do servidor físico, peça à docente -- muitas instalações Proxmox expõem métricas por SNMP para fins de demonstração.

\section{Introdução}
Este guião coloca em prática a arquitetura SNMP estudada na aula teórica: uma VM desempenha o papel de agente (com o daemon \texttt{snmpd} a expor a sua MIB), e outra o papel de manager, consultando essa MIB através de \texttt{snmpget} e \texttt{snmpwalk}. Configura-se depois o envio de uma trap, para observar o mecanismo assíncrono de notificação, e por fim um utilizador SNMPv3, para comparar a segurança das duas abordagens.

Começamos com SNMPv2c e uma \textit{community string}, por ser a configuração mais simples. Na Parte 3 veremos, capturando tráfego real, porque é que essa simplicidade tem um custo -- e como o SNMPv3 o resolve.

\section{Partes 1 e 2 -- Agente SNMPv2c, consultas e traps}
\begin{enumerate}[label=\arabic*.]
  \item \textbf{Instalar o agente na VM gerida.} Via SSH à VM, instale o pacote \texttt{snmpd}:
\begin{lstlisting}
$ sudo apt update
$ sudo apt install snmpd -y
\end{lstlisting}

  \item \textbf{Configurar a community string.} Faça uma cópia de segurança do ficheiro de configuração original e edite-o:
\begin{lstlisting}
$ sudo cp /etc/snmp/snmpd.conf /etc/snmp/snmpd.conf.bak
$ sudo nano /etc/snmp/snmpd.conf
\end{lstlisting}
  Substitua as linhas de restrição de acesso (\texttt{rocommunity}) por uma linha que defina a sua própria community string, com acesso de leitura a partir da rede local, por exemplo:
\begin{lstlisting}
rocommunity icnturma 192.168.1.0/24
\end{lstlisting}
  Ajuste o intervalo de rede ao usado no seu laboratório.

  \item \textbf{Permitir acesso à MIB-II completa (opcional, ambiente de aprendizagem).} Por omissão, algumas distribuições restringem o agente a um subconjunto reduzido da MIB. Para expor a árvore de sistema e de interfaces de rede, garanta que a linha seguinte está presente (ou remova restrições de \texttt{view} demasiado apertadas):
\begin{lstlisting}
view   systemonly  included   .1.3.6.1.2.1
\end{lstlisting}

  \item \textbf{Reiniciar o serviço.} Aplique a nova configuração:
\begin{lstlisting}
$ sudo systemctl restart snmpd
$ sudo systemctl enable snmpd
$ sudo systemctl status snmpd
\end{lstlisting}
  Confirme que o serviço fica ativo (\texttt{active (running)}).

  \item \textbf{Instalar as ferramentas de cliente na máquina manager.} Na segunda máquina (fora da VM gerida):
\begin{lstlisting}
$ sudo apt update
$ sudo apt install snmp snmp-mibs-downloader -y
\end{lstlisting}

  \item \textbf{Consultar o uptime do sistema com \texttt{snmpget}.} O OID \texttt{1.3.6.1.2.1.1.3.0} corresponde ao \texttt{sysUpTime} standard da MIB-II. Substitua \texttt{IP\_da\_VM} pelo endereço da VM gerida:
\begin{lstlisting}
$ snmpget -v2c -c icnturma IP_da_VM 1.3.6.1.2.1.1.3.0
\end{lstlisting}
  A resposta deve indicar há quanto tempo o agente está em execução.

  \item \textbf{Percorrer a tabela de interfaces de rede com \texttt{snmpwalk}.} A tabela \texttt{ifTable} da MIB-II (ramo \texttt{1.3.6.1.2.1.2}) descreve todas as interfaces de rede conhecidas pelo agente:
\begin{lstlisting}
$ snmpwalk -v2c -c icnturma IP_da_VM 1.3.6.1.2.1.2.2.1.2
\end{lstlisting}
  Este comando lista os nomes (\texttt{ifDescr}) de todas as interfaces de rede da VM. Experimente também percorrer a árvore de sistema completa:
\begin{lstlisting}
$ snmpwalk -v2c -c icnturma IP_da_VM 1.3.6.1.2.1.1
\end{lstlisting}

  \item \textbf{Traduzir OIDs em nomes legíveis.} Compare a saída numérica com a saída traduzida pelos ficheiros MIB:
\begin{lstlisting}
$ snmpwalk -v2c -c icnturma -On IP_da_VM 1.3.6.1.2.1.1
$ snmpwalk -v2c -c icnturma     IP_da_VM 1.3.6.1.2.1.1
\end{lstlisting}
  A opção \texttt{-On} força a apresentação numérica. Se os nomes não aparecerem, ative o carregamento das MIBs:
\begin{lstlisting}
$ export MIBS=ALL
\end{lstlisting}

  \item \textbf{Observar tipos de dados: Counter vs Gauge.} Consulte os contadores de bytes recebidos numa interface, duas vezes, com alguns segundos de intervalo:
\begin{lstlisting}
$ snmpget -v2c -c icnturma IP_da_VM 1.3.6.1.2.1.2.2.1.10.2
$ sleep 10
$ snmpget -v2c -c icnturma IP_da_VM 1.3.6.1.2.1.2.2.1.10.2
\end{lstlisting}
  O OID \texttt{ifInOctets} é do tipo \texttt{Counter32}: só cresce. Registe os dois valores e calcule a diferença dividida pelo intervalo -- obtém assim o débito médio em bytes por segundo. Este é exatamente o cálculo que qualquer sistema de monitorização faz internamente.

  Compare com um valor do tipo \texttt{Gauge}, que sobe e desce livremente:
\begin{lstlisting}
$ snmpwalk -v2c -c icnturma IP_da_VM 1.3.6.1.2.1.25.2.3.1
\end{lstlisting}

  \item \textbf{Configurar o recetor de traps.} Na máquina manager, instale e configure um recetor de traps simples:
\begin{lstlisting}
$ sudo apt install snmptrapd -y
$ sudo nano /etc/snmp/snmptrapd.conf
\end{lstlisting}
  Adicione a linha seguinte, com a mesma community string usada no agente:
\begin{lstlisting}
authCommunity log,execute,net icnturma
\end{lstlisting}
  Reinicie o serviço:
\begin{lstlisting}
$ sudo systemctl restart snmptrapd
\end{lstlisting}

  \item \textbf{Configurar a VM gerida para enviar traps.} De volta à VM gerida, edite novamente o \texttt{snmpd.conf} e acrescente o destino da trap (substitua \texttt{IP\_do\_manager} pelo endereço da máquina manager):
\begin{lstlisting}
trap2sink IP_do_manager icnturma
\end{lstlisting}
  Reinicie o \texttt{snmpd} para aplicar a alteração:
\begin{lstlisting}
$ sudo systemctl restart snmpd
\end{lstlisting}

  \item \textbf{Disparar uma trap de teste.} Ainda na VM gerida, envie manualmente uma trap de teste para o manager:
\begin{lstlisting}
$ snmptrap -v2c -c icnturma IP_do_manager '' 1.3.6.1.6.3.1.1.5.3
\end{lstlisting}
  Na máquina manager, observe o registo do \texttt{snmptrapd} (por omissão em \texttt{/var/log/syslog} ou na consola, consoante a configuração) e confirme a receção da notificação.
\end{enumerate}

\section{Parte 3 -- SNMPv3 e comparação de segurança}
\begin{enumerate}[label=\arabic*.]

  \item \textbf{Capturar tráfego SNMPv2c.} Na VM agente, inicie uma captura na porta 161:
\begin{lstlisting}
$ sudo apt install -y tcpdump
$ sudo tcpdump -i any -A -n port 161 -c 20
\end{lstlisting}
  A partir do manager, faça uma consulta:
\begin{lstlisting}
$ snmpget -v2c -c icnturma IP_da_VM 1.3.6.1.2.1.1.5.0
\end{lstlisting}
  Procure na captura a \textit{community string} \texttt{icnturma}. Confirme que está visível em texto simples -- é esta a fragilidade discutida na aula teórica.

  \item \textbf{Criar um utilizador SNMPv3.} Pare o agente antes de o criar:
\begin{lstlisting}
$ sudo systemctl stop snmpd
$ sudo net-snmp-create-v3-user -ro -A "AuthPass2026" -a SHA \
    -X "PrivPass2026" -x AES aluno_icn
$ sudo systemctl start snmpd
\end{lstlisting}
  O comando cria um utilizador de leitura com autenticação SHA e cifra AES -- o nível \texttt{authPriv}.

  \item \textbf{Consultar com SNMPv3} a partir do manager:
\begin{lstlisting}
$ snmpget -v3 -l authPriv -u aluno_icn \
    -a SHA -A "AuthPass2026" \
    -x AES -X "PrivPass2026" \
    IP_da_VM 1.3.6.1.2.1.1.5.0
\end{lstlisting}

  \item \textbf{Repetir a captura, agora com SNMPv3:}
\begin{lstlisting}
$ sudo tcpdump -i any -A -n port 161 -c 20
\end{lstlisting}
  Repita a consulta v3 e procure na captura o nome do utilizador e o valor devolvido. Registe o que consegue e o que não consegue ler.

  \item \textbf{Testar os níveis de segurança.} Compare o comportamento com autenticação apenas (\texttt{authNoPriv}) e verifique se os dados passam a ser legíveis na captura:
\begin{lstlisting}
$ snmpget -v3 -l authNoPriv -u aluno_icn -a SHA -A "AuthPass2026" \
    IP_da_VM 1.3.6.1.2.1.1.5.0
\end{lstlisting}

  \item \textbf{Testar o controlo de acesso.} Tente uma escrita com o utilizador de leitura e registe a resposta obtida:
\begin{lstlisting}
$ snmpset -v3 -l authPriv -u aluno_icn -a SHA -A "AuthPass2026" \
    -x AES -X "PrivPass2026" \
    IP_da_VM 1.3.6.1.2.1.1.6.0 s "Nova localizacao"
\end{lstlisting}

\end{enumerate}

\section{Entregável / Questões de Verificação}
\begin{enumerate}[label=\alph*)]
  \setlength\itemsep{0.2cm}
  \item Capture o resultado dos comandos \texttt{snmpget} (uptime) e \texttt{snmpwalk} (interfaces de rede) e submeta-os, juntamente com o registo da trap recebida no manager.
  \item Que diferença existe entre uma operação \texttt{GET} e uma \texttt{TRAP}? Qual delas é iniciada pelo agente? Relacione com a distinção entre \emph{polling} e \emph{push}.
  \item Submeta as duas capturas de \texttt{tcpdump} (SNMPv2c e SNMPv3, com anonimização das passwords). O que era legível em cada caso? Que nível de segurança do SNMPv3 usou?
  \item Porque motivo o SNMPv2c com community strings é considerado inseguro para uso em redes não confiáveis? Que mecanismos do SNMPv3 resolvem o problema?
  \item Registe os dois valores de \texttt{ifInOctets} e o intervalo entre as leituras. Calcule o débito médio em bytes por segundo. Porque é que o valor absoluto de um \texttt{Counter32} não é, por si só, informativo?
  \item Qual a diferença entre um tipo \texttt{Counter} e um tipo \texttt{Gauge}? Dê um exemplo de métrica adequada a cada um.
  \item O que representa o OID \texttt{1.3.6.1.2.1.1.3.0}? A que ramo da MIB pertence e o que significa cada um dos seus níveis?
  \item O que aconteceu quando tentou o \texttt{snmpset} com um utilizador de leitura? Que mecanismo do SNMPv3 impediu a operação?
  \item Se tivesse de monitorizar as VMs desta turma, que três métricas escolheria para gerar alertas, e com que limiares? Justifique tendo em conta o problema da fadiga de alertas.
\end{enumerate}

\end{document}
